% 和文 2 段組（学会誌の体裁）。xelatex で組む。
%
% japanese-onecolumn.tex との違いは**欧文の術語の書き方**:
%   こちらは前後に空白を置かない（\texttt{attention}機構）。xeCJK や LuaTeX-ja を
%   使う和文原稿ではこちらが普通で、紙面にも空白は出ない。
%   したがって和文と欧文の境目で折り返しても空白を入れてはいけない。
%   空白を置く流儀（japanese-onecolumn.tex）とは正解が逆になるので、
%   「継ぎ目に必ず空白を入れる」でも「継ぎ目に絶対入れない」でも、どちらかが落ちる。
%
% 英文要旨を付けてあるのは、同じ文書の中で欧文の折り返し（空白が要る）と
% 和文の折り返し（空白が入ってはいけない）が両方起きることを測るため。
%
% 正解の取り出し方は scripts/real-papers/*.tex と同じ（本文 1 文 = \s{...} を 1 行）。
% ----------------------------------------------------------------------------
% 計測（npm run measure）への足し方
%
%   1. この .tex と .pdf と .truth.json を scripts/real-papers/ に置く。
%      measureRealPapers() は real-papers/*.pdf を総当たりで読むので、
%      measure-quality.ts 側に足す記述は要らない。
%   2. 組み直しには **xelatex** が要る（和文の書体を使うため pdflatex では通らない）。
%      scripts/build-real-papers.ts は pdflatex 決め打ちなので、
%      「\usepackage{fontspec} があれば xelatex」のような分岐を足す。
%      正解（.truth.json）の作り方は欧文と同じ scripts/real-paper-truth.ts のまま
%      で通る（和文の \s{...} もそのまま拾える）。
%   3. verifyInlineMath（pdftotext との突き合わせ）は \m{...}{...} を 1 つも
%      使っていないので、この原稿では何もしない。
% ----------------------------------------------------------------------------
\documentclass[10pt,a4paper,twocolumn]{article}

\usepackage{fontspec}
\setmainfont{Noto Serif CJK JP}
\setmonofont{DejaVu Sans Mono}
\XeTeXlinebreaklocale "ja"
% 和文の字間に入れる伸縮。pTeX の \kanjiskip（10pt で plus/minus 0.4pt）に合わせる。
% ここを大きく取ると、pdf.js が字間の空きを空白文字と見なして
% （TRACKING_SPACE_FACTOR = 0.102 em）文字片の中に空白を撒く。実在の和文組版は
% そこまで空けないので、現実的な版面にしてから測る。
\XeTeXlinebreakskip=0pt plus 0.4pt minus 0.4pt
\usepackage[margin=2.0cm]{geometry}
% 柱間は **LaTeX の既定（\columnsep = 10pt）のまま**にする。
%
% 和文の学会誌でも既定のまま出す原稿は珍しくないし、なにより
% 「8mm に広げれば通る」標本では、和文で段組の検出が落ちることを測れない。
% 和文は 1 文字がほぼ 1em 進むので、段の切れ目の下限を「字送りの 1.5 倍」で
% 決めると 10pt 本文で 14.9pt になり、10pt の柱間がそのまま下回る
% （欧文は字送り 4.4pt 前後なので同じ 10pt が 2.3 倍あって通る）。
\usepackage{fancyhdr}

\newcommand{\s}[1]{#1}
% 参考文献の見出しも和文にする。和文誌の体裁ではこちらが普通。
\renewcommand{\refname}{参考文献}

\pagestyle{fancy}
\fancyhf{}
\fancyhead[L]{情報処理学会論文誌}
\fancyfoot[C]{\thepage}

\title{和文と欧文が隣り合う版面からの文の切り出し}
\author{坂比賀 悠真\\ 山手大学 情報科学研究科}
\date{}

\begin{document}
\maketitle

\begin{abstract}
\s{This paper studies sentence segmentation for papers whose body text is written in Japanese.}
\s{Existing tools assume that words are separated by spaces, which never happens inside a Japanese paragraph.}
\s{We report the error rate of a rule based segmenter on three Japanese corpora.}%
\end{abstract}

\section{はじめに}

\s{和文の学会誌では、本文を日本語で書き、術語だけを欧文のまま残すことが多い。}%
\s{たとえば\texttt{attention}機構や\texttt{transformer}系の模型はそのまま表記される。}%
\s{この書き方では、和文と欧文の境目に空白が入らない。}%
\s{空白を手掛かりに語を区切る解析器は、この版面で必ず誤る。}%
\s{さらに和文の行は語の途中でも折り返すため、行を繋ぐ位置の判断が難しい。}%

\s{本稿では、文字の種類から継ぎ目の空白の要否を決める規則を示す。}%
\s{規則は決定的であり、辞書も学習データも必要としない。}%
\s{三つの和文コーパスで評価し、既存の実装と比較した。}%

\section{関連研究}

\s{英語の文境界推定では終止符と大文字が主な手掛かりであり、\texttt{nltk}の実装が広く使われている。}%
\s{和文では句点がほぼ唯一の手掛かりであり、括弧の中の句点をどう扱うかが問題になる \cite{sato2021}。}%
\s{形態素解析器を前段に置く手法もあるが、辞書の整備に費用がかかる \cite{tanaka2019}。}%
\s{本稿が対象とするのは、辞書を用意できない場面である。}%

\section{提案手法}

\s{継ぎ目の両側が漢字・ひらがな・カタカナ・全角記号のいずれかであれば、空白を入れない。}%
\s{両側が欧文であれば、これまでどおり空白を入れる。}%
\s{和文と欧文の境目では、同じ文書の行の途中に現れる境目の書き方に合わせる。}%
\s{この文書のように空白を置かない流儀であれば、折り返しでも空白を入れない。}%
\s{規則の適用順序は固定であり、同じ入力には必ず同じ出力を返す。}%

\section{実装}

\s{実装は\texttt{typescript}で書かれており、外部の依存を持たない。}%
\s{入力は版面から取り出した文字片の列で、出力は文の列である。}%
\s{文字片は座標を持つので、行にまとめる処理を先に通してから文に割る。}%
\s{行にまとめる処理は段の切れ目を空白の帯として検出し、左の段から順に読む。}%
\s{この順序を誤ると左右の段が一行の中で混ざり、文がすべて壊れる。}%

\s{和文の版面では、段の切れ目の判定に使う字送りが欧文の二倍になる。}%
\s{一文字がほぼ一エムを占めるため、同じ幅の柱間でも相対的には狭く見える。}%
\s{この性質を考えずに閾値を決めると、和文の二段組で段組の検出そのものが失敗する。}%
\s{実装ではページごとに字送りを測り、その倍数で閾値を決めている。}%

\section{評価}

\s{三つのコーパスは情報処理学会論文誌、人工知能学会誌、および和文要旨付きのプレプリントから構成される。}%
\s{いずれも人手で文の境界を付与し、二人の作業者の一致率は九十七パーセントであった。}%
\s{提案手法の再現率は三つのコーパスでそれぞれ九十九、九十八、九十七であった。}%
\s{既存の実装は、折り返しのたびに語の途中へ空白を入れるため、再現率が三割に留まった。}%
\s{誤りの多くは、箇条書きと表の見出しが本文に混ざる場合に集中していた。}%

\s{処理時間は百ページの論文で一秒未満であり、実用上の問題は無い。}%
\s{記憶領域は文字片の数に比例し、百ページの論文で数十メガバイトに収まる。}%

\section{考察}

\s{残った誤りの半分は、見出しが本文として扱われたことに起因する。}%
\s{和文の見出しは番号と本文の間に空白が一つ入るだけで、欧文のような大文字の手掛かりが無い。}%
\s{文字の大きさだけを手掛かりにすると、本文と同じ大きさで組まれた小見出しを取り逃がす。}%
\s{この問題は文の切り出しとは独立であり、行の役割を決める処理の側で解くべきである。}%

\s{もう一つの誤りは、柱と本文が同じ段に並ぶ場合に起きる。}%
\s{柱は複数のページに同じ字面で現れるので、ページ数が少ない文書では判別が難しい。}%
\s{この場合も文の切り出しではなく、行の役割を決める処理の問題である。}%

\section{誤りの分析}

\s{誤りを種類ごとに数えると、見出しに起因するものが最も多く、次に柱に起因するものが続いた。}%
\s{見出しに起因する誤りは、見出しと直後の本文が一枚の札に同居する形で現れる。}%
\s{柱に起因する誤りは、ページの変わり目で本文の途中に柱の文字列が挟まる形で現れる。}%
\s{どちらも読む側からは「原文がそう書いてある」ようにしか見えないため、気付きにくい。}%
\s{気付きにくい誤りは、気付ける誤りよりも害が大きいと我々は考える。}%

\s{三つ目の種類は、図表の見出しが本文として扱われる誤りである。}%
\s{和文の図表の見出しは「図 1」「表 2」の形をとり、欧文の版面とは接頭辞が異なる。}%
\s{接頭辞の一覧に和文の語を加えれば解けるが、雑誌ごとに表記が揺れるため万全ではない。}%
\s{我々は、接頭辞に頼らず、版面の位置と前後の行との間隔から判断する方法を検討している。}%

\s{四つ目の種類は、箇条書きの各項目が一つの文として扱われる誤りである。}%
\s{箇条書きの項目は句点で終わらないことが多く、次の項目と繋がってしまう。}%
\s{行頭の記号を手掛かりにすれば切れるが、記号の種類は版面によって異なる。}%
\s{この種類の誤りは全体の一割に満たないため、優先度は低いと判断した。}%

\section{他の言語への拡張}

\s{中国語は日本語と同じく分かち書きをしないので、同じ規則がそのまま使える。}%
\s{一方、朝鮮語は分かち書きをするため、行の継ぎ目には欧文と同じく空白が要る。}%
\s{同じ東アジアの文字であっても、空白の扱いは言語によって異なる。}%
\s{したがって、文字の範囲を「東アジアの文字」とひとまとめにする実装は誤りである。}%
\s{我々の実装は、漢字・ひらがな・カタカナ・全角記号のみを対象とし、ハングルを含めない。}%

\s{縦書きの版面については、行の向きが変わるだけで規則そのものは変わらない。}%
\s{ただし、版面から行を組み立てる処理は縦書きを想定していないため、別に手当てが要る。}%
\s{振り仮名は本文と重なる位置に置かれるため、行の組み立ての段階で落とす必要がある。}%
\s{これらは今後の課題である。}%

\section{実運用での知見}

\s{本手法を社内の読解支援の道具に組み込み、半年間の運用で得た知見を述べる。}%
\s{利用者は和文誌の論文と欧文誌の論文を同じ道具で読むため、言語ごとに設定を切り替える余地は無い。}%
\s{したがって、言語の判別は文書の中身から自動で行う必要がある。}%
\s{我々は文書全体ではなく、継ぎ目ごとにその両側の文字だけを見る方式を採った。}%
\s{この方式であれば、一つの文書に和文と欧文の段落が混在していても正しく扱える。}%

\s{運用の初期には、和文の要旨だけが崩れるという報告が繰り返し寄せられた。}%
\s{原因は、要旨が本文よりも狭い幅で組まれ、折り返しの回数が本文より多かったことである。}%
\s{折り返しのたびに語が割れるため、幅が狭いほど被害が大きくなる。}%
\s{この事実は、崩れの量が版面の幅に依存することを示している。}%

\s{利用者からの報告で最も多かったのは、訳文ではなく原文が読めないという苦情であった。}%
\s{対訳の道具であるから、原文が崩れていれば訳文の正しさは意味を持たない。}%
\s{原文の忠実さは、対訳の道具にとって最も優先すべき性質であると我々は考える。}%

\s{運用の後半では、崩れの報告は月に一件未満まで減った。}%
\s{残った報告はいずれも見出しと柱に起因するもので、本文の連結に起因するものは無かった。}%

\section{実装上の注意}

\s{文字片の継ぎ目に空白を入れるかどうかは、隙間の大きさだけでは決められない。}%
\s{和文は一文字がほぼ一エム進むため、句読点の詰めだけで欧文の空白一つ分を超える。}%
\s{我々の実装では、継ぎ目の両側が和文であるときに限って閾値を引き上げている。}%
\s{本物の空白は、多くの場合は抽出器が空白文字として返してくるため、隙間から推測する必要が無い。}%

\s{抽出器そのものが空白を作ってしまう場合もある。}%
\s{字送りが一定の割合を超えると、版面に空白が無くても空白文字が差し込まれる。}%
\s{両端揃えの和文は行の伸縮を字間に配るため、狭い段ではこの条件に触れやすい。}%
\s{この場合、一行のほぼすべての字間に空白が入るので、行全体を見れば見分けられる。}%
\s{著者名や参考文献の姓名にある本物の空白は、一行に一つか二つしか現れない。}%
\s{我々は、個数と割合の両方が閾値を超えた行に限って、和文どうしの間の空白を落としている。}%

\s{行末のハイフンにも注意が必要である。}%
\s{和文は語の途中で折り返すときに記号を足さないため、行末のハイフンは欧文にしか現れない。}%
\s{和文の書体で欧文を組むと、行末のハイフンが通常の記号とは異なる符号で取り出されることがある。}%
\s{この場合、符号の違いを吸収しないと、割れた語が画面にそのまま残る。}%

\s{以上の三つは、いずれも版面を見ただけでは気付けない種類の崩れである。}%
\s{崩れの有無を数値で測る仕組みを先に用意したことが、これらの発見に繋がった。}%
\s{測る仕組みが無ければ、和文の論文が読めないという報告を受けても、原因の切り分けができない。}%
\s{我々は、和文の版面を計測の標本に加えることを強く勧める。}%

\section{おわりに}

\s{和文の版面から文を切り出すための、文字の種類にもとづく規則を示した。}%
\s{三つのコーパスで既存の実装を大きく上回る結果が得られた。}%
\s{今後は、縦書きと振り仮名を含む版面への拡張を検討したい。}%

\begin{thebibliography}{9}
% 参考文献も和文で書く。和文誌では引用先も和文であることが普通で、
% 「[1] の後ろが欧文の大文字」しか見ない判定では拾えない。
\bibitem{sato2021}
佐藤 花子, 鈴木 一郎. 日本語文書における文境界推定の再検討. 情報処理学会論文誌, 62(3):512--523, 2021.
\bibitem{tanaka2019}
田中 明. 学術文書の構造解析. 人工知能学会誌, 34(2):88--97, 2019.
\end{thebibliography}

\end{document}
